\documentclass[11pt]{article}
\usepackage[margin=25mm]{geometry}
\usepackage{amsmath,amssymb,amsthm}
\usepackage[T1]{fontenc}
\usepackage[colorlinks=true,urlcolor=blue]{hyperref}
\newtheorem{theorem}{Theorem}
\setlength{\parindent}{0pt}\setlength{\parskip}{5pt}
\begin{document}
\begin{center}{\Large A stationary nilsequence with correlation identically one}\\[5pt]{\large Conjecture 00000001924}\\[5pt]zbw123456 \quad 14 October 2026\end{center}

Let $G=(\mathbb R,+)$ and $\Gamma=\mathbb Z$, so $G/\Gamma=\mathbb R/\mathbb Z$ is a compact one-step nilmanifold. Take the constant polynomial orbit $g(n)=0$ and the continuous, Lipschitz observable $F\equiv1$. Then $a_n=F(g(n)+\Gamma)=1$ is a basic polynomial nilsequence. It is invariant under every integer shift, so it satisfies stationarity without any extra condition on the shift.
\begin{theorem}
The normalized correlation of this nilsequence with the zero-frequency linear phase equals one for every $N\ge1$ and admits no $O(N^{-1/2})$ bound.
\end{theorem}
\begin{proof}
For $e(\alpha n)=\exp(2\pi i\alpha n)$, set $\alpha=0$. Then
\[
\frac1N\sum_{n=0}^{N-1} a_n\overline{e(0n)}=1.
\]
If the asserted decay held, some real constant $C$ would satisfy $1\le C N^{-1/2}$ for all sufficiently large $N$. The right side tends to zero, giving $1\le0$, a contradiction.
\end{proof}
\textbf{Hypotheses.} The official statement does not require a mean-zero observable, a nonresonance condition, a nonzero frequency, or subtraction of the mean. Such restrictions would change the conjecture. The witness uses the ordinary normalized correlation; starting the sum at one gives the same value.

\textbf{Lean correspondence.} \texttt{circleSequence} uses Mathlib's \texttt{AddCircle 1}, its actual additive rotation, and a bundled continuous observable. \texttt{witness\_one} and \texttt{stationary} prove the claimed sequence and shift invariance. \texttt{correlation} uses complex exponentials, conjugation and a finite sum; \texttt{correlation\_one} evaluates it for every nonzero $N$. \texttt{no\_weyl\_bound} rules out every eventual square-root bound using a proved limit.

\textbf{Definition source.} Green and Tao, \href{https://arxiv.org/abs/0807.1736}{The M\"obius function is strongly orthogonal to nilsequences}, and Leibman, \href{https://arxiv.org/abs/1205.4004}{Nilsequences, null-sequences, and multiple correlation sequences}. The present problem concerns linear phases, not the M\"obius function.


\textbf{Verification.} Lean 4.19.0 with Mathlib pinned to
\texttt{c44e0c8ee63ca166450922a373c7409c5d26b00b}.
The accompanying project and verification transcript give the exact theorem and axiom audit. No proof placeholders, native evaluator proof bypass, or custom axiom declarations are used. This is a prepared submission; official acceptance has not been claimed.
\end{document}
